%@ pages 328
%@ status translated
% Chapter I, section 1: Grothendieck universes (p. 328).

\chapter{A Few Recalls on Categories}
\markboth{Chapter I. A Few Recalls on Categories}{}

The purpose of this chapter is to recall a few definitions and notations
concerning categories. We suppose that the reader is acquainted with
Chapters~I and~II of [10]. However, since time has done its work since 1957,
we add a few supplements to Grothendieck's article; these supplements are due
in large part to Grothendieck himself, and they will be expounded in a
forthcoming work [7]. This is why we allow ourselves to omit most of the
proofs.

\section{Grothendieck universes}
--- A set $\Univ$ is a \emph{universe} if the following axioms hold:

\begin{description}
\item[$\mathrm{U}_1$] If $(X_i)_{i\in I}$ is a family of sets belonging to
$\Univ$, and if $I$ is an element of $\Univ$, then the union $\bigcup_i X_i$
is an element of $\Univ$.

\item[$\mathrm{U}_2$] If $x$ belongs to $\Univ$, then the set $\{x\}$ is an
element belonging to $\Univ$.

\item[$\mathrm{U}_3$] If $x$ belongs to $X$ and if $X$ belongs to $\Univ$,
then $x$ belongs to $\Univ$.

\item[$\mathrm{U}_4$] If $X$ is a set belonging to $\Univ$, the set
$\Powerset(X)$ of parts of $X$ is an element of $\Univ$.

\item[$\mathrm{U}_5$] The pair $(x,y)$ is an element of $\Univ$ if and only
if $x$ and $y$ are elements of $\Univ$.
\end{description}

It is convenient to add to the usual axioms of set theory an axiom ensuring
that every set belongs to a universe. In that case there exists a smallest
universe containing a given set. We choose once and for all a universe $\Univ$
which will not ``vary'' throughout what follows. Of course, we suppose $\Univ$
large enough that the set $\mathbf{Z}$ of integers, and possibly other sets,
are elements of $\Univ$. The reader may exercise himself in proving the
following corollaries of the axioms:

\begin{itemize}
\item If $Y$ is contained in $X$, and if $X$ belongs to $\Univ$, then $Y$
      belongs to $\Univ$.
\item If $X$ and $Y$ are two sets belonging to $\Univ$, then the union
      $X \cup Y$ and the Cartesian product $X \times Y$ are elements of
      $\Univ$.
\item If $x$ and $y$ belong to $\Univ$, then the set $\{x,y\}$ belongs to
      $\Univ$.
\item If $(X_i)_{i\in I}$ is a family of sets belonging to $\Univ$, and if
      $I$ belongs to $\Univ$, then the product $\prod_{i\in I} X_i$ belongs
      to $\Univ$.
\item If $X$ is a set belonging to $\Univ$, then the cardinality of $X$ is
      strictly smaller than the cardinality of $\Univ$.
\end{itemize}